\section{Introduction: The Intersection of Arbitrary-Precision Numerics and Causal Deep Learning}

The integration of high-fidelity numerical analysis with deep learning represents a paradigm shift in the study of chaotic dynamical systems. The specific research stack under analysis—comprising the Brutus arbitrary-precision integrator written in C++, a Python-based data processing pipeline, and a Causal Training Labeled Physics-Informed Neural Network (CTL-PINN)—targets one of the most enduring challenges in classical mechanics: the Three-Body Problem. This system is characterized by non-linear interactions, extreme sensitivity to initial conditions, and phase space structures defined by the coexistence of stable invariant tori and chaotic seas. The successful modeling of such a system using neural networks requires a rigorous reassessment of what constitutes ``high-quality'' training data. Unlike traditional supervised learning tasks where data quality is often synonymous with label accuracy or noise reduction, the chaotic N-body problem imposes strict geometric and dynamical constraints—specifically symplectic conservation and phase space continuity—that must be preserved to ensure the learned model is physically viable.

The primary objective of this report is to define a comprehensive strategy for assessing training data quality prior to full model training. The analysis is driven by the unique requirements of the CTL-PINN architecture, which enforces temporal causality in the learning process, thereby making the model highly sensitive to the temporal coherence and physical consistency of the input data. We evaluate the comparative efficacy of static pre-training validation versus dynamic trial training, ultimately proposing a hybrid ``Staged Gate'' protocol. Furthermore, we develop a set of mathematically grounded metrics tailored to the arbitrary-precision nature of the Brutus dataset, ensuring that the inductive biases of the neural network align with the symplectic geometry of the underlying Hamiltonian physics.

\subsection{The Computational Landscape: Brutus and the Bulirsch-Stoer Algorithm}

To understand the data quality requirements, one must first analyze the source of the data. The Brutus code differs fundamentally from the symplectic integrators (such as Leapfrog or Wisdom-Holman) typically used in large-scale N-body simulations. Symplectic integrators preserve the geometric properties of the Hamiltonian flow (such as phase space volume) over long time scales but do so at the cost of local accuracy, often introducing a bounded energy error that oscillates but does not grow secularly. In contrast, Brutus utilizes the Bulirsch-Stoer algorithm combined with arbitrary-precision arithmetic to generate ``converged'' solutions.

The Bulirsch-Stoer method is an extrapolation scheme. It computes the solution at a future time $t + H$ by subdividing the interval $H$ into substeps, performing the integration with a modified midpoint method, and then extrapolating the results to zero step size using rational functions. Brutus iteratively increases the precision (word length $L$) and decreases the tolerance parameter $\epsilon$ until the solution converges to a specified number of significant digits. This implies that a Brutus-generated trajectory is not merely an approximation; within the specified precision, it is a ``true'' numerical solution that shadows the exact mathematical trajectory of the system.

However, this high fidelity comes with a specific data characteristic: non-uniform temporal sampling. The adaptive step-size controller in Brutus shrinks the time step $\Delta t$ dramatically during close encounters (pericenter passages) to maintain accuracy as the gravitational forces $F \propto 1/r^2$ diverge. Conversely, during excursive phases where bodies are widely separated, $\Delta t$ increases. While this is efficient for integration, it presents a ``spectral bias'' challenge for neural networks, which typically assume uniform sampling or at least sampling densities that correlate with the local Lipschitz constant of the target function in a way that is compatible with gradient descent dynamics. Consequently, ``good'' data for Brutus—defined by convergence—is not necessarily ``good'' data for a CTL-PINN without rigorous validation of its sampling density and local stiffness properties.

\subsection{The Learning Architecture: Causal Training Labeled PINN (CTL-PINN)}

The target model is a CTL-PINN. This architecture represents a refinement of the standard PINN formulation designed specifically for time-dependent evolution equations. Standard PINNs often suffer from ``temporal causality violation,'' where the optimizer attempts to minimize the residual at a late time $t=T$ before having adequately resolved the solution at $t=0$. In chaotic systems, this is catastrophic because the solution at $t=T$ is exponentially sensitive to the trajectory at $t < T$. If the early trajectory is not learned correctly, the physics loss at later times provides incorrect gradient information, leading to convergence to unphysical local minima.

Causal training addresses this by introducing a weighting mechanism $w(t)$ into the loss function:
\begin{equation}
    \mathcal{L}(\theta) = \sum_{i=1}^{N_t} w_i \mathcal{L}_i(\theta)
\end{equation}
where the weights $w_i$ are dynamically adjusted based on the cumulative loss at previous time steps. Typically, $w_i \propto \exp(-\lambda \sum_{j < i} \mathcal{L}_j)$. This forces the network to learn the trajectory sequentially, respecting the arrow of time.

This architectural choice imposes strict requirements on the training data. The data must be \textit{causally continuous}. A gap in the data, a region of unphysical noise, or a ``stiffness spike'' (an under-resolved close encounter) at $t_{early}$ will result in a high residual $\mathcal{L}_{early}$. The causal weighting mechanism will then suppress the weights for all $t > t_{early}$ to zero, effectively halting the training process. Thus, for CTL-PINNs, data quality is not just about global accuracy; it is about the \textit{continuity of learnability} along the temporal axis. A single ``bad'' segment renders the entire subsequent trajectory useless for training.